\documentclass[11pt,a4paper]{article}
\usepackage[margin=21mm]{geometry}
\usepackage{amsmath,amssymb,booktabs,array}
\setlength{\parindent}{0pt}
\setlength{\parskip}{5pt}
\pagestyle{plain}
\begin{document}
\begin{center}{\Large\bfseries Homework 5: Concise Answers}\end{center}

\section*{A (p.\ 1): Proof in words}
With the bar denoting negation, the statement is
\[
[(a\lor b)\land(\neg a\lor c)]\to\neg(b\lor c).
\]
\textbf{The statement is false as written.} Suppose $a$ and $c$ are true
and $b$ is false. The first clause is true because $a$ is true, and the
second clause is true because $c$ is true. Thus the premise is true.
However, $b\lor c$ is true, so its negation is false.
This gives a true premise and a false conclusion, making the implication
false. Therefore it cannot be proved as a valid implication.
\[
\boxed{a=T,\quad b=F,\quad c=T\quad\text{is a counterexample}.}
\]
The valid resolution conclusion would be $b\lor c$.

\section*{B (p.\ 1, second clause set): Resolution}
Number the given clauses $C_1,\ldots,C_6$.
Each subsequent row is the resolvent of the indicated two clauses.
\begin{center}
\renewcommand{\arraystretch}{1.16}
\begin{tabular}{c|l|l}
\toprule
Clause&Expression&Reason\\
\midrule
$C_1$&$p\lor q\lor\neg r\lor s$&Given\\
$C_2$&$\neg p\lor r\lor s$&Given\\
$C_3$&$\neg q\lor\neg r$&Given\\
$C_4$&$p\lor\neg s$&Given\\
$C_5$&$\neg p\lor\neg r$&Given\\
$C_6$&$r$&Given\\
\midrule
$C_7$&$\neg q$&$C_3,C_6$; resolve on $r$\\
$C_8$&$\neg p$&$C_5,C_6$; resolve on $r$\\
$C_9$&$\neg s$&$C_4,C_8$; resolve on $p$\\
$C_{10}$&$p\lor q\lor s$&$C_1,C_6$; resolve on $r$\\
$C_{11}$&$p\lor s$&$C_{10},C_7$; resolve on $q$\\
$C_{12}$&$s$&$C_{11},C_8$; resolve on $p$\\
$C_{13}$&$\square$ (empty clause)&$C_{12},C_9$; resolve on $s$\\
\bottomrule
\end{tabular}
\end{center}
The empty clause is false under every assignment. Since resolution is
sound, the original clause set is $\boxed{\text{unsatisfiable}}$.
The contradiction already follows from $C_1,C_3,C_4,C_5,C_6$;
$C_2$ is not needed.

\newpage
\section*{C (p.\ 4, Exercise 58): Simultaneously true disjunctions}
Let
\[
\begin{aligned}
D_1&=p\lor\neg q,&D_2&=\neg p\lor q,&D_3&=q\lor r,\\
D_4&=q\lor\neg r,&D_5&=\neg q\lor\neg r.
\end{aligned}
\]
Choose $p=T,\ q=T,\ r=F$. Then
\[
\begin{aligned}
D_1&=T\lor F=T, &D_2&=F\lor T=T,\\
D_3&=T\lor F=T, &D_4&=T\lor T=T,\\
D_5&=F\lor T=T.
\end{aligned}
\]
All five are simultaneously true. Since there are only five disjunctions,
\[
\boxed{\text{The maximum is }5.}
\]

\section*{D (p.\ 5): Equations and implications}
\textbf{Equations.} Over the real numbers,
\[
x=y,\quad y=z,\quad z=-x
\quad\Longrightarrow\quad x=y=z=-x
\quad\Longrightarrow\quad 2x=0.
\]
Hence $x=y=z=0$. Conversely, this assignment satisfies all three equations.
\[
\boxed{\text{Yes, exactly when }x=y=z=0.}
\]

\textbf{Propositional logic.} Let
\[
\Phi=(A\to B)\land(B\to C)\land(C\to\neg A).
\]
An implication is false only when its antecedent is true and its consequent
is false.
\begin{center}
\renewcommand{\arraystretch}{1.16}
\begin{tabular}{ccc|ccc|c}
\toprule
$A$&$B$&$C$&$A\to B$&$B\to C$&$C\to\neg A$&$\Phi$\\
\midrule
T&T&T&T&T&F&F\\
T&T&F&T&F&T&F\\
T&F&T&F&T&F&F\\
T&F&F&F&T&T&F\\
F&T&T&T&T&T&T\\
F&T&F&T&F&T&F\\
F&F&T&T&T&T&T\\
F&F&F&T&T&T&T\\
\bottomrule
\end{tabular}
\end{center}
\[
\boxed{\text{Choice (b): true simultaneously in some cases, but not all.}}
\]
Exactly three assignments work: $(A,B,C)=(F,T,T),(F,F,T),(F,F,F)$.
Equivalently, $\Phi\equiv\neg A\land(\neg B\lor C)$, which is satisfiable
but is not a tautology.
\end{document}
